%& -shell-escape
% !TeX root = thesis-abstract.tex
% !TeX encoding = UTF-8
% !TeX program = XeLaTeX

\documentclass{.local/lib/classes/ndvr-super-article-standalone}

\begin{document}
    \section{Основные результаты}

    На защиту выносятся:
    \begin{itemize}
        \item 
            метод сравнения и~поиска нечетких дубликатов видео, 
            который с~помощью разбиения видео на события 
            сводит задачу к сравнению и~поиску временных рядов \cite{
                Nikitin:CMASS:2013,
                Nikitin:CMASS:2015,
                Nikitin:HSE:2013,
                Nikitin:MGPPU:2013,
                Nikitin:NC:2015,
                Nikitin:NGU:2013,
                Nikitin:RSDN:2014,
                Nikitin:SYTT:2013
            }; 
        \item
            теорема о нечетких дубликатах видео \cite{Lukin:ALG:2017};
        \item
            способ описания процесса обработки видео
            с помощью декларативного языка фильтров \cite{Lukin:ALG:2017}, 
            пример приведён ранее~— на рисунке \ref{fig:mean-sign-diff-event-filter};
        \item 
            грамматика декларативного языка фильтров \cite{Lukin:ALG:2017};
        \item
            алгоритм сегментации видео
            на основе сравнения скользящих 
            средних оценок сигнала \cite{Lukin:ALG:2017}.
    \end{itemize}

    % \begin{figure}[H]
    %     \includegraphics[width=0.9\linewidth]
    %     {.local/vec/out/pdf/ffmpeg-filter-code.pdf}
    %     \caption{
    %          Фильтр FFMpeg 
    %     }
    %     \label{fig:mean-sign-diff-event-filter}
    % \end{figure}

\end{document}


